%% ============================================================
%% INTREPIDCORE — ic-editorial.tex (guide de rédaction v7.0)
%% Bibliothèque du style éditorial abstrait et géométrique.
%%   1. Jetons de couleur des deux ambiances (Marque et Terre)
%%   2. Motifs de l'Atlas (ic-motifs.tex) et cartes de motif
%%   3. Ouvertures de partie (ambiance Marque) et couverture éditoriale
%%   4. Bibliothèque de l'ambiance Terre (trait irrégulier)
%% Prérequis : xelatex, tikz (calc, arrows.meta, decorations.pathmorphing),
%% xcolor, les couleurs papier / intrepid / accent / cyanIC / vertLab du
%% préambule, et la macro \papertexture.
%% ============================================================

%% ── 1. Jetons de couleur ────────────────────────────────────
%% Ambiance Marque : alignée sur la palette de production
\colorlet{mnavy}{intrepid}
\colorlet{morange}{accent}
\colorlet{mcyan}{cyanIC}
\colorlet{mvert}{vertLab}
\colorlet{mpapier}{papier}
\definecolor{mblanc}{HTML}{FFFDF9}
\definecolor{mencre}{HTML}{1C2329}
\definecolor{mgris}{HTML}{65717C}
\definecolor{mligne}{HTML}{D8D8D1}
\definecolor{mpoint}{HTML}{DFDBD4}
%% Ambiance Terre
\definecolor{tcotta}{RGB}{196,98,61}
\definecolor{tcottaL}{RGB}{219,154,116}
\definecolor{creme}{RGB}{243,234,217}
\definecolor{beige}{RGB}{227,212,184}
\definecolor{sable}{RGB}{217,183,126}
\definecolor{argile}{RGB}{233,205,185}
\definecolor{noirenc}{RGB}{30,27,24}

%% ── 2. Motifs et cartes ─────────────────────────────────────
\newcommand{\Motif}[1]{\csname Motif#1\endcsname}
%% ============================================================
%% INTREPIDCORE — ic-motifs.tex
%% Les 24 motifs de l'Atlas graphique IntrepidCore (familles A à D),
%% dessinés en TikZ natif. Source : Atlas_IntrepidCore_24_motifs_source_LaTeX.zip,
%% fichier atlas_intrepidcore.tex (octobre 2026).
%% Adaptation : les noms de couleur sont préfixés par « m » (mnavy, morange,
%% mcyan, mvert, mblanc, mencre, mgris, mligne, mpoint, mpapier) pour ne pas
%% redéfinir les couleurs xcolor « orange » et « cyan ». Les couleurs mnavy,
%% morange, mcyan, mvert, mpapier sont alignées sur la palette du guide
%% (définies dans ic-editorial.tex).
%% Cadre de dessin : 150 x 100 unités ; appel : \Motif{A01} dans une
%% tikzpicture [x=1pt,y=1pt,scale=...]. Statut : proposition graphique.
%% ============================================================
%% A01 | Noyau et satellites
\expandafter\def\csname MotifA01\endcsname{%
\begin{scope}
\begin{scope}[shift={(0,0)}]
\begin{scope}[xscale=1,yscale=1]
\foreach \dx in {8,23,...,143} {\foreach \dy in {7,22,...,97} {\fill[mpoint] (\dx,\dy) circle[radius=0.65];}}
\draw[draw=mnavy,fill=none,line width=1.6pt] (25,70) -- (76,49);
\draw[draw=mnavy,fill=none,line width=1.6pt] (36,20) -- (76,49);
\draw[draw=mnavy,fill=none,line width=1.6pt] (125,69) -- (76,49);
\draw[draw=mnavy,fill=none,line width=1.6pt] (118,22) -- (76,49);
\path[draw=none,fill=mnavy] (75,49) circle[radius=23];
\path[draw=none,fill=mblanc] (75,49) circle[radius=12];
\path[draw=mnavy,fill=mblanc,line width=1.8pt,rounded corners=2pt] (15,60) rectangle (35,80);
\path[draw=mnavy,fill=mblanc,line width=1.8pt,rounded corners=2pt] (28,12) rectangle (44,28);
\path[draw=mnavy,fill=mblanc,line width=1.8pt,rounded corners=2pt] (115,59) rectangle (135,79);
\path[draw=mnavy,fill=mblanc,line width=1.8pt,rounded corners=2pt] (110,14) rectangle (126,30);
\path[draw=none,fill=morange] (76,49) circle[radius=3];
\draw[draw=morange,fill=none,line width=1.5pt] (35,70) -- (116,20);
\end{scope}
\end{scope}
\end{scope}
}

%% A02 | Strates emboîtées
\expandafter\def\csname MotifA02\endcsname{%
\begin{scope}
\begin{scope}[shift={(0,0)}]
\begin{scope}[xscale=1,yscale=1]
\foreach \dx in {8,23,...,143} {\foreach \dy in {7,22,...,97} {\fill[mpoint] (\dx,\dy) circle[radius=0.65];}}
\path[draw=mnavy,fill=mblanc,line width=2.2pt,rounded corners=2pt] (25,18) rectangle (100,80);
\path[draw=mnavy,fill=mblanc,line width=2.2pt,rounded corners=2pt] (43,32) rectangle (118,94);
\path[draw=none,fill=mnavy,rounded corners=2pt] (56,13) rectangle (118,76);
\path[draw=none,fill=mblanc] (88,47) circle[radius=19];
\path[draw=none,fill=morange] (88,47) circle[radius=9];
\path[draw=morange,fill=mblanc,line width=2pt,rounded corners=2pt] (33,59) rectangle (53,70);
\draw[draw=morange,fill=none,line width=1.1pt] (28,13) -- (123,85);
\end{scope}
\end{scope}
\end{scope}
}

%% A03 | Maillage modulaire
\expandafter\def\csname MotifA03\endcsname{%
\begin{scope}
\begin{scope}[shift={(0,0)}]
\begin{scope}[xscale=1,yscale=1]
\foreach \dx in {8,23,...,143} {\foreach \dy in {7,22,...,97} {\fill[mpoint] (\dx,\dy) circle[radius=0.65];}}
\path[draw=mnavy,fill=mnavy,line width=1.4pt,rounded corners=1.5pt] (20,17) rectangle (44,38);
\path[draw=none,fill=morange,rounded corners=1.5pt] (47,17) rectangle (71,38);
\path[draw=mnavy,fill=mblanc,line width=1.4pt,rounded corners=1.5pt] (74,17) rectangle (98,38);
\path[draw=mnavy,fill=mnavy,line width=1.4pt,rounded corners=1.5pt] (101,17) rectangle (125,38);
\path[draw=mnavy,fill=mblanc,line width=1.4pt,rounded corners=1.5pt] (33,42) rectangle (57,63);
\path[draw=mnavy,fill=mblanc,line width=1.4pt,rounded corners=1.5pt] (60,42) rectangle (84,63);
\path[draw=none,fill=morange,rounded corners=1.5pt] (87,42) rectangle (111,63);
\path[draw=mnavy,fill=mblanc,line width=1.4pt,rounded corners=1.5pt] (114,42) rectangle (138,63);
\path[draw=none,fill=morange,rounded corners=1.5pt] (20,67) rectangle (44,88);
\path[draw=mnavy,fill=mnavy,line width=1.4pt,rounded corners=1.5pt] (47,67) rectangle (71,88);
\path[draw=mnavy,fill=mblanc,line width=1.4pt,rounded corners=1.5pt] (74,67) rectangle (98,88);
\path[draw=mnavy,fill=mblanc,line width=1.4pt,rounded corners=1.5pt] (101,67) rectangle (125,88);
\draw[draw=morange,fill=none,line width=0.9pt] (42,29) -- (56,42);
\draw[draw=morange,fill=none,line width=0.9pt] (69,29) -- (83,42);
\draw[draw=morange,fill=none,line width=0.9pt] (96,29) -- (110,42);
\end{scope}
\end{scope}
\end{scope}
}

%% A04 | Équilibre des contraintes
\expandafter\def\csname MotifA04\endcsname{%
\begin{scope}
\begin{scope}[shift={(0,0)}]
\begin{scope}[xscale=1,yscale=1]
\foreach \dx in {8,23,...,143} {\foreach \dy in {7,22,...,97} {\fill[mpoint] (\dx,\dy) circle[radius=0.65];}}
\draw[draw=mnavy,fill=none,line width=1.65pt] (42,24) -- (109,22);
\draw[draw=mnavy,fill=none,line width=1.65pt] (42,24) -- (75,85);
\draw[draw=mnavy,fill=none,line width=1.65pt] (109,22) -- (75,85);
\draw[draw=morange,fill=none,line width=1.65pt] (42,24) -- (76,52);
\draw[draw=morange,fill=none,line width=1.65pt] (109,22) -- (76,52);
\draw[draw=morange,fill=none,line width=1.65pt] (75,85) -- (76,52);
\path[draw=mnavy,fill=mblanc,line width=2.2pt,rounded corners=3pt] (31,13) rectangle (53,35);
\path[draw=mnavy,fill=mblanc,line width=2.2pt,rounded corners=3pt] (98,11) rectangle (120,33);
\path[draw=mnavy,fill=mblanc,line width=2.2pt,rounded corners=3pt] (64,74) rectangle (86,96);
\path[draw=none,fill=mnavy] (76,52) circle[radius=15];
\path[draw=none,fill=mblanc] (76,52) circle[radius=6];
\path[draw=none,fill=morange] (42,24) circle[radius=3];
\end{scope}
\end{scope}
\end{scope}
}

%% A05 | Interface commune
\expandafter\def\csname MotifA05\endcsname{%
\begin{scope}
\begin{scope}[shift={(0,0)}]
\begin{scope}[xscale=1,yscale=1]
\foreach \dx in {8,23,...,143} {\foreach \dy in {7,22,...,97} {\fill[mpoint] (\dx,\dy) circle[radius=0.65];}}
\path[draw=none,fill=mnavy,rounded corners=4pt] (17,27) rectangle (79,75);
\path[draw=mnavy,fill=mblanc,line width=2.6pt,rounded corners=4pt] (71,20) rectangle (132,68);
\path[draw=none,fill=mblanc] (75,47) circle[radius=19];
\path[draw=none,fill=mcyan] (75,47) circle[radius=11];
\draw[draw=mblanc,fill=none,line width=2pt] (31,34) -- (45,34);
\draw[draw=mblanc,fill=none,line width=2pt] (31,48) -- (45,48);
\draw[draw=mblanc,fill=none,line width=2pt] (31,62) -- (45,62);
\draw[draw=morange,fill=none,line width=2pt] (104,31) -- (116,31);
\draw[draw=morange,fill=none,line width=2pt] (104,47) -- (116,47);
\draw[draw=morange,fill=none,line width=2pt] (104,63) -- (116,63);
\draw[draw=morange,fill=none,line width=2pt] (50,47) -- (63,47);
\path[draw=none,fill=morange] (63,47) -- (59.2366,48.8646) -- (59.2366,45.1354) -- cycle;
\end{scope}
\end{scope}
\end{scope}
}

%% A06 | Ossature assemblée
\expandafter\def\csname MotifA06\endcsname{%
\begin{scope}
\begin{scope}[shift={(0,0)}]
\begin{scope}[xscale=1,yscale=1]
\foreach \dx in {8,23,...,143} {\foreach \dy in {7,22,...,97} {\fill[mpoint] (\dx,\dy) circle[radius=0.65];}}
\path[draw=none,fill=mnavy] (26,18) rectangle (51,49);
\path[draw=none,fill=mnavy] (99,18) rectangle (124,49);
\path[draw=mnavy,fill=mblanc,line width=1.7pt] (40,62) circle[radius=10];
\path[draw=mnavy,fill=mblanc,line width=1.7pt] (54,78) circle[radius=10];
\path[draw=mnavy,fill=mblanc,line width=1.7pt] (75,83) circle[radius=10];
\path[draw=mnavy,fill=mblanc,line width=1.7pt] (96,78) circle[radius=10];
\path[draw=mnavy,fill=mblanc,line width=1.7pt] (111,62) circle[radius=10];
\draw[draw=morange,fill=none,line width=3pt] (117,46) arc[start angle=0,end angle=180,x radius=42,y radius=42];
\draw[draw=mnavy,fill=none,line width=2pt] (106,46) arc[start angle=0,end angle=180,x radius=31,y radius=31];
\path[draw=mnavy,fill=mblanc,line width=2.2pt] (57,16) rectangle (93,46);
\path[draw=none,fill=morange] (75,32) circle[radius=7];
\end{scope}
\end{scope}
\end{scope}
}

%% B01 | Chaîne de transformations
\expandafter\def\csname MotifB01\endcsname{%
\begin{scope}
\begin{scope}[shift={(0,0)}]
\begin{scope}[xscale=1,yscale=1]
\foreach \dx in {8,23,...,143} {\foreach \dy in {7,22,...,97} {\fill[mpoint] (\dx,\dy) circle[radius=0.65];}}
\path[draw=none,fill=mnavy,rounded corners=2pt] (10,38) rectangle (34,62);
\path[draw=mcyan,fill=mblanc,line width=2pt] (65,50) circle[radius=13];
\path[draw=none,fill=morange] (94,37) -- (108,64) -- (122,37) -- cycle;
\path[draw=none,fill=mnavy] (141,50) circle[radius=8];
\draw[draw=mnavy,fill=none,line width=2pt] (34,50) -- (52,50);
\path[draw=none,fill=mnavy] (52,50) -- (48.2366,51.8646) -- (48.2366,48.1354) -- cycle;
\draw[draw=mnavy,fill=none,line width=2pt] (78,50) -- (92,50);
\path[draw=none,fill=mnavy] (92,50) -- (88.2366,51.8646) -- (88.2366,48.1354) -- cycle;
\draw[draw=mnavy,fill=none,line width=2pt] (123,50) -- (132,50);
\path[draw=none,fill=mnavy] (132,50) -- (128.2366,51.8646) -- (128.2366,48.1354) -- cycle;
\path[draw=none,fill=mblanc] (22,50) circle[radius=5];
\path[draw=none,fill=mcyan] (65,50) circle[radius=4];
\end{scope}
\end{scope}
\end{scope}
}

%% B02 | Relais successifs
\expandafter\def\csname MotifB02\endcsname{%
\begin{scope}
\begin{scope}[shift={(0,0)}]
\begin{scope}[xscale=1,yscale=1]
\foreach \dx in {8,23,...,143} {\foreach \dy in {7,22,...,97} {\fill[mpoint] (\dx,\dy) circle[radius=0.65];}}
\path[draw=mnavy,fill=none,line width=2pt] (11,26) .. controls (40,10) and (53,40) .. (78,49) .. controls (95,63) and (106,77) .. (139,78);
\path[draw=mnavy,fill=mblanc,line width=1.9pt] (16,26) circle[radius=7];
\path[draw=none,fill=mcyan] (16,26) circle[radius=3.5];
\path[draw=mnavy,fill=mblanc,line width=1.9pt] (48,29) circle[radius=9];
\path[draw=none,fill=mcyan] (48,29) circle[radius=4.5];
\path[draw=mnavy,fill=mblanc,line width=1.9pt] (80,49) circle[radius=11];
\path[draw=none,fill=mcyan] (80,49) circle[radius=5.5];
\path[draw=mnavy,fill=mblanc,line width=1.9pt] (113,72) circle[radius=14];
\path[draw=none,fill=morange] (113,72) circle[radius=7];
\draw[draw=morange,fill=none,line width=1.6pt] (122,77) -- (142,82);
\path[draw=none,fill=morange] (142,82) -- (137.8967,82.8961) -- (138.8012,79.2783) -- cycle;
\end{scope}
\end{scope}
\end{scope}
}

%% B03 | Branches réunifiées
\expandafter\def\csname MotifB03\endcsname{%
\begin{scope}
\begin{scope}[shift={(0,0)}]
\begin{scope}[xscale=1,yscale=1]
\foreach \dx in {8,23,...,143} {\foreach \dy in {7,22,...,97} {\fill[mpoint] (\dx,\dy) circle[radius=0.65];}}
\path[draw=none,fill=mnavy,rounded corners=2pt] (8,41) rectangle (27,60);
\path[draw=none,fill=morange] (136,50) circle[radius=12];
\draw[draw=mnavy,fill=none,line width=2pt] (27,50) -- (43,50);
\path[draw=mcyan,fill=none,line width=2.15pt] (43,50) .. controls (58,50) and (59,83) .. (86,79) .. controls (109,77) and (107,50) .. (124,50);
\path[draw=mnavy,fill=none,line width=2.15pt] (43,50) .. controls (58,50) and (82,50) .. (124,50);
\path[draw=morange,fill=none,line width=2.15pt] (43,50) .. controls (56,48) and (60,18) .. (86,18) .. controls (109,18) and (111,50) .. (124,50);
\path[draw=mnavy,fill=mblanc,line width=1pt] (43,50) circle[radius=3];
\draw[draw=mnavy,fill=none,line width=1pt] (124,50) -- (126,50);
\path[draw=none,fill=mnavy] (126,50) -- (122.2366,51.8646) -- (122.2366,48.1354) -- cycle;
\end{scope}
\end{scope}
\end{scope}
}

%% B04 | Filtre et sélection
\expandafter\def\csname MotifB04\endcsname{%
\begin{scope}
\begin{scope}[shift={(0,0)}]
\begin{scope}[xscale=1,yscale=1]
\foreach \dx in {8,23,...,143} {\foreach \dy in {7,22,...,97} {\fill[mpoint] (\dx,\dy) circle[radius=0.65];}}
\path[draw=mnavy,fill=mblanc,line width=1.4pt] (12,27) circle[radius=4];
\path[draw=mnavy,fill=mblanc,line width=1.4pt] (23,42) circle[radius=4];
\path[draw=mnavy,fill=mblanc,line width=1.4pt] (35,61) circle[radius=4];
\path[draw=mnavy,fill=mblanc,line width=1.4pt] (20,77) circle[radius=4];
\path[draw=mnavy,fill=mblanc,line width=1.4pt] (39,25) circle[radius=4];
\path[draw=mnavy,fill=mblanc,line width=1.4pt] (50,74) circle[radius=4];
\path[draw=mnavy,fill=mblanc,line width=1.4pt] (54,45) circle[radius=4];
\path[draw=mnavy,fill=mblanc,line width=1.4pt] (46,58) circle[radius=4];
\draw[draw=mnavy,fill=none,line width=2.2pt] (67,21) -- (67,79);
\draw[draw=mcyan,fill=none,line width=2pt] (63,24) -- (69,24);
\draw[draw=mcyan,fill=none,line width=2pt] (63,30) -- (69,30);
\draw[draw=mcyan,fill=none,line width=2pt] (63,36) -- (69,36);
\draw[draw=mcyan,fill=none,line width=2pt] (63,42) -- (69,42);
\draw[draw=mcyan,fill=none,line width=2pt] (63,48) -- (69,48);
\draw[draw=mcyan,fill=none,line width=2pt] (63,54) -- (69,54);
\draw[draw=mcyan,fill=none,line width=2pt] (63,60) -- (69,60);
\draw[draw=mcyan,fill=none,line width=2pt] (63,66) -- (69,66);
\draw[draw=mcyan,fill=none,line width=2pt] (63,72) -- (69,72);
\draw[draw=mcyan,fill=none,line width=2pt] (63,78) -- (69,78);
\path[draw=none,fill=mnavy,rounded corners=2pt] (85,43) rectangle (97,57);
\path[draw=none,fill=morange,rounded corners=2pt] (101,43) rectangle (113,57);
\path[draw=none,fill=morange,rounded corners=2pt] (117,43) rectangle (129,57);
\path[draw=none,fill=morange,rounded corners=2pt] (135,43) rectangle (147,57);
\draw[draw=morange,fill=none,line width=1.6pt] (73,50) -- (79,50);
\path[draw=none,fill=morange] (79,50) -- (75.2366,51.8646) -- (75.2366,48.1354) -- cycle;
\end{scope}
\end{scope}
\end{scope}
}

%% B05 | Boucle de rétroaction
\expandafter\def\csname MotifB05\endcsname{%
\begin{scope}
\begin{scope}[shift={(0,0)}]
\begin{scope}[xscale=1,yscale=1]
\foreach \dx in {8,23,...,143} {\foreach \dy in {7,22,...,97} {\fill[mpoint] (\dx,\dy) circle[radius=0.65];}}
\path[draw=mnavy,fill=mblanc,line width=1.9pt] (75,84) circle[radius=13];
\path[draw=none,fill=mcyan] (75,84) circle[radius=5];
\path[draw=mnavy,fill=mblanc,line width=1.9pt] (111,50) circle[radius=13];
\path[draw=none,fill=morange] (111,50) circle[radius=5];
\path[draw=mnavy,fill=mblanc,line width=1.9pt] (75,16) circle[radius=13];
\path[draw=none,fill=mnavy] (75,16) circle[radius=5];
\path[draw=mnavy,fill=mblanc,line width=1.9pt] (39,50) circle[radius=13];
\path[draw=none,fill=mcyan] (39,50) circle[radius=5];
\draw[draw=mnavy,fill=none,line width=2pt] (107.5144,59.9406) arc[start angle=17,end angle=73,x radius=34,y radius=34];
\path[draw=none,fill=mnavy] (84.9406,82.5144) -- (87.849,79.7682) -- (88.8874,83.1646) -- cycle;
\draw[draw=mnavy,fill=none,line width=2pt] (65.0594,82.5144) arc[start angle=107,end angle=163,x radius=34,y radius=34];
\path[draw=none,fill=mnavy] (42.4856,59.9406) -- (45.2318,62.849) -- (41.8354,63.8874) -- cycle;
\draw[draw=morange,fill=none,line width=2pt] (42.4856,40.0594) arc[start angle=197,end angle=253,x radius=34,y radius=34];
\path[draw=none,fill=morange] (65.0594,17.4856) -- (62.151,20.2318) -- (61.1126,16.8354) -- cycle;
\draw[draw=mnavy,fill=none,line width=2pt] (84.9406,17.4856) arc[start angle=287,end angle=343,x radius=34,y radius=34];
\path[draw=none,fill=mnavy] (107.5144,40.0594) -- (104.7682,37.151) -- (108.1646,36.1126) -- cycle;
\end{scope}
\end{scope}
\end{scope}
}

%% B06 | Distribution ramifiée
\expandafter\def\csname MotifB06\endcsname{%
\begin{scope}
\begin{scope}[shift={(0,0)}]
\begin{scope}[xscale=1,yscale=1]
\foreach \dx in {8,23,...,143} {\foreach \dy in {7,22,...,97} {\fill[mpoint] (\dx,\dy) circle[radius=0.65];}}
\path[draw=none,fill=mnavy,rounded corners=2pt] (14,36) rectangle (39,65);
\path[draw=none,fill=mblanc] (26,50) circle[radius=7];
\path[draw=mnavy,fill=none,line width=1.6pt] (39,50) .. controls (74,50) and (79,82) .. (121,82);
\path[draw=mcyan,fill=mblanc,line width=1.8pt] (121,82) circle[radius=9];
\path[draw=mnavy,fill=none,line width=1.6pt] (39,50) .. controls (74,50) and (79,50) .. (127,50);
\path[draw=mcyan,fill=mblanc,line width=1.8pt] (127,50) circle[radius=9];
\path[draw=mnavy,fill=none,line width=1.6pt] (39,50) .. controls (74,50) and (79,18) .. (121,18);
\path[draw=mcyan,fill=mblanc,line width=1.8pt] (121,18) circle[radius=9];
\path[draw=none,fill=morange] (95,82) circle[radius=2];
\path[draw=none,fill=morange] (101,50) circle[radius=2];
\path[draw=none,fill=morange] (95,18) circle[radius=2];
\draw[draw=mcyan,fill=none,line width=0.75pt] (132.3562,40.0414) arc[start angle=-50,end angle=50,x radius=13,y radius=13];
\draw[draw=mcyan,fill=none,line width=0.75pt] (136.213,35.4452) arc[start angle=-50,end angle=50,x radius=19,y radius=19];
\end{scope}
\end{scope}
\end{scope}
}

%% C01 | Dispersion autour d’un repère
\expandafter\def\csname MotifC01\endcsname{%
\begin{scope}
\begin{scope}[shift={(0,0)}]
\begin{scope}[xscale=1,yscale=1]
\foreach \dx in {8,23,...,143} {\foreach \dy in {7,22,...,97} {\fill[mpoint] (\dx,\dy) circle[radius=0.65];}}
\path[draw=mvert,fill=none,line width=1.4pt] (75,50) circle[radius=36];
\path[draw=none,fill=mvert] (49,43) circle[radius=5];
\path[draw=none,fill=mvert] (60,68) circle[radius=3.6];
\path[draw=none,fill=mvert] (70,48) circle[radius=3.6];
\path[draw=none,fill=morange] (83,60) circle[radius=5];
\path[draw=none,fill=mvert] (94,36) circle[radius=3.6];
\path[draw=none,fill=mvert] (48,67) circle[radius=3.6];
\path[draw=none,fill=mvert] (99,72) circle[radius=5];
\path[draw=none,fill=mvert] (80,30) circle[radius=3.6];
\path[draw=none,fill=morange] (65,25) circle[radius=3.6];
\path[draw=none,fill=mvert] (104,48) circle[radius=5];
\path[draw=none,fill=mvert] (42,25) circle[radius=3.6];
\path[draw=none,fill=mvert] (76,75) circle[radius=3.6];
\path[draw=none,fill=mvert] (37,55) circle[radius=5];
\path[draw=none,fill=mvert] (113,61) circle[radius=3.6];
\path[draw=mnavy,fill=mblanc,line width=2pt] (75,50) circle[radius=6];
\draw[draw=mnavy,fill=none,line width=1pt] (64,50) -- (86,50);
\draw[draw=mnavy,fill=none,line width=1pt] (75,39) -- (75,61);
\end{scope}
\end{scope}
\end{scope}
}

%% C02 | Famille de scénarios
\expandafter\def\csname MotifC02\endcsname{%
\begin{scope}
\begin{scope}[shift={(0,0)}]
\begin{scope}[xscale=1,yscale=1]
\foreach \dx in {8,23,...,143} {\foreach \dy in {7,22,...,97} {\fill[mpoint] (\dx,\dy) circle[radius=0.65];}}
\path[draw=mvert,fill=mblanc,line width=1.9pt] (20,18) circle[radius=6.7];
\path[draw=none,fill=mvert,rounded corners=1.3pt] (34.5,13.5) rectangle (47.5,26.5);
\path[draw=mvert,fill=mblanc,line width=1.9pt] (62,22) circle[radius=6.7];
\path[draw=none,fill=mvert,rounded corners=1.3pt] (76.5,17.5) rectangle (89.5,30.5);
\path[draw=mvert,fill=mblanc,line width=1.9pt] (104,19) circle[radius=6.7];
\path[draw=none,fill=mvert,rounded corners=1.3pt] (118.5,14.5) rectangle (131.5,27.5);
\path[draw=none,fill=mvert,rounded corners=1.3pt] (17.5,35.5) rectangle (30.5,48.5);
\path[draw=mvert,fill=mblanc,line width=1.9pt] (45,44) circle[radius=6.7];
\path[draw=none,fill=mvert,rounded corners=1.3pt] (59.5,32.5) rectangle (72.5,45.5);
\path[draw=mvert,fill=mblanc,line width=1.9pt] (87,41) circle[radius=6.7];
\path[draw=none,fill=mvert,rounded corners=1.3pt] (101.5,36.5) rectangle (114.5,49.5);
\path[draw=mvert,fill=mblanc,line width=1.9pt] (129,38) circle[radius=6.7];
\path[draw=mvert,fill=mblanc,line width=1.9pt] (21,59) circle[radius=6.7];
\path[draw=none,fill=mvert,rounded corners=1.3pt] (35.5,54.5) rectangle (48.5,67.5);
\path[draw=mvert,fill=mblanc,line width=1.9pt] (63,63) circle[radius=6.7];
\path[draw=none,fill=morange] (84,58) circle[radius=9];
\path[draw=mvert,fill=mblanc,line width=1.9pt] (105,60) circle[radius=6.7];
\path[draw=none,fill=mvert,rounded corners=1.3pt] (119.5,55.5) rectangle (132.5,68.5);
\path[draw=none,fill=mvert,rounded corners=1.3pt] (18.5,76.5) rectangle (31.5,89.5);
\path[draw=mvert,fill=mblanc,line width=1.9pt] (46,78) circle[radius=6.7];
\path[draw=none,fill=mvert,rounded corners=1.3pt] (60.5,73.5) rectangle (73.5,86.5);
\path[draw=mvert,fill=mblanc,line width=1.9pt] (88,82) circle[radius=6.7];
\path[draw=none,fill=mvert,rounded corners=1.3pt] (102.5,77.5) rectangle (115.5,90.5);
\path[draw=mvert,fill=mblanc,line width=1.9pt] (130,79) circle[radius=6.7];
\path[draw=mnavy,fill=none,line width=0.85pt,rounded corners=2pt] (16,12) rectangle (144,92);
\end{scope}
\end{scope}
\end{scope}
}

%% C03 | Centre mobile
\expandafter\def\csname MotifC03\endcsname{%
\begin{scope}
\begin{scope}[shift={(0,0)}]
\begin{scope}[xscale=1,yscale=1]
\foreach \dx in {8,23,...,143} {\foreach \dy in {7,22,...,97} {\fill[mpoint] (\dx,\dy) circle[radius=0.65];}}
\path[draw=mvert,fill=none,line width=1.8pt] (76,50) circle[radius=41];
\path[draw=mnavy,fill=none,line width=1.8pt] (79,51.5) circle[radius=34];
\path[draw=mvert,fill=none,line width=1.8pt] (82,53) circle[radius=27];
\path[draw=mnavy,fill=none,line width=1.8pt] (85,54.5) circle[radius=20];
\path[draw=mvert,fill=none,line width=1.8pt] (88,56) circle[radius=13];
\path[draw=none,fill=morange] (89,55) circle[radius=7];
\draw[draw=morange,fill=none,line width=1pt] (33,18) -- (40,12);
\draw[draw=morange,fill=none,line width=1pt] (41,18) -- (48,12);
\draw[draw=morange,fill=none,line width=1pt] (49,18) -- (56,12);
\draw[draw=morange,fill=none,line width=1pt] (57,18) -- (64,12);
\draw[draw=morange,fill=none,line width=1pt] (65,18) -- (72,12);
\end{scope}
\end{scope}
\end{scope}
}

%% C04 | Domaines superposés
\expandafter\def\csname MotifC04\endcsname{%
\begin{scope}
\begin{scope}[shift={(0,0)}]
\begin{scope}[xscale=1,yscale=1]
\foreach \dx in {8,23,...,143} {\foreach \dy in {7,22,...,97} {\fill[mpoint] (\dx,\dy) circle[radius=0.65];}}
\path[draw=mvert,fill=none,line width=2pt] (32,55) .. controls (32,68.5) and (46.5,73) .. (61,73) .. controls (84.2,73) and (90,67.6) .. (90,55) .. controls (90,38.8) and (78.4,37) .. (61,37) .. controls (37.8,37) and (32,46) .. (32,55) -- cycle;
\path[draw=mnavy,fill=none,line width=2pt] (54,56) .. controls (54,73.25) and (71,79) .. (88,79) .. controls (115.2,79) and (122,72.1) .. (122,56) .. controls (122,35.3) and (108.4,33) .. (88,33) .. controls (60.8,33) and (54,44.5) .. (54,56) -- cycle;
\path[draw=mvert,fill=none,line width=2pt] (40,38) .. controls (40,50.75) and (60,55) .. (80,55) .. controls (112,55) and (120,49.9) .. (120,38) .. controls (120,22.7) and (104,21) .. (80,21) .. controls (48,21) and (40,29.5) .. (40,38) -- cycle;
\path[draw=mnavy,fill=none,line width=2pt] (23,38) .. controls (23,53) and (38,58) .. (53,58) .. controls (77,58) and (83,52) .. (83,38) .. controls (83,20) and (71,18) .. (53,18) .. controls (29,18) and (23,28) .. (23,38) -- cycle;
\path[draw=none,fill=morange] (75,47) circle[radius=5];
\end{scope}
\end{scope}
\end{scope}
}

%% C05 | Perturbation locale
\expandafter\def\csname MotifC05\endcsname{%
\begin{scope}
\begin{scope}[shift={(0,0)}]
\begin{scope}[xscale=1,yscale=1]
\foreach \dx in {8,23,...,143} {\foreach \dy in {7,22,...,97} {\fill[mpoint] (\dx,\dy) circle[radius=0.65];}}
\path[draw=none,fill=mvert] (15,25) circle[radius=4];
\path[draw=none,fill=mnavy] (15,54.2074) circle[radius=4];
\path[draw=none,fill=mvert] (15,79.5465) circle[radius=4];
\path[draw=none,fill=mvert] (33,28.2211) circle[radius=4];
\path[draw=none,fill=mnavy] (33,54.9583) circle[radius=4];
\path[draw=none,fill=mvert] (33,77.1369) circle[radius=4];
\path[draw=none,fill=mvert] (51,29.9272) circle[radius=4];
\path[draw=none,fill=mnavy] (51,53.3773) circle[radius=4];
\path[draw=none,fill=mvert] (51,73.7223) circle[radius=4];
\path[draw=none,fill=mvert] (69,29.316) circle[radius=4];
\path[draw=none,fill=mnavy] (69,50.2079) circle[radius=4];
\path[draw=none,fill=mvert] (69,70.9086) circle[radius=4];
\path[draw=none,fill=mvert] (87,26.6749) circle[radius=4];
\path[draw=none,fill=mnavy] (87,46.9407) circle[radius=4];
\path[draw=none,fill=mvert] (87,70.0192) circle[radius=4];
\path[draw=none,fill=mvert] (105,23.2461) circle[radius=6];
\path[draw=none,fill=morange] (105,45.1123) circle[radius=6];
\path[draw=none,fill=mvert] (105,71.4723) circle[radius=6];
\path[draw=none,fill=mvert] (123,20.6421) circle[radius=4];
\path[draw=none,fill=mnavy] (123,45.5827) circle[radius=4];
\path[draw=none,fill=mvert] (123,74.5846) circle[radius=4];
\path[draw=none,fill=mvert] (141,20.0877) circle[radius=4];
\path[draw=none,fill=mnavy] (141,48.1306) circle[radius=4];
\path[draw=none,fill=mvert] (141,77.8922) circle[radius=4];
\draw[draw=mvert,fill=none,line width=1pt] (95,70) arc[start angle=90,end angle=180,x radius=20,y radius=20];
\draw[draw=mvert,fill=none,line width=1pt] (95,77) arc[start angle=90,end angle=180,x radius=27,y radius=27];
\draw[draw=mvert,fill=none,line width=1pt] (95,84) arc[start angle=90,end angle=180,x radius=34,y radius=34];
\draw[draw=mvert,fill=none,line width=1pt] (95,91) arc[start angle=90,end angle=180,x radius=41,y radius=41];
\end{scope}
\end{scope}
\end{scope}
}

%% C06 | Éventail de possibilités
\expandafter\def\csname MotifC06\endcsname{%
\begin{scope}
\begin{scope}[shift={(0,0)}]
\begin{scope}[xscale=1,yscale=1]
\foreach \dx in {8,23,...,143} {\foreach \dy in {7,22,...,97} {\fill[mpoint] (\dx,\dy) circle[radius=0.65];}}
\draw[draw=mnavy,fill=none,line width=1.5pt] (29,29) -- (126.8148,49.7912);
\path[draw=none,fill=mvert] (126.8148,49.7912) circle[radius=3];
\draw[draw=mvert,fill=none,line width=1.5pt] (29,29) -- (135.8436,65.7892);
\draw[draw=mvert,fill=none,line width=1.5pt] (29,29) -- (126.0698,76.3441);
\path[draw=none,fill=mvert] (126.0698,76.3441) circle[radius=3];
\draw[draw=mnavy,fill=none,line width=1.5pt] (29,29) -- (115.3831,85.0978);
\draw[draw=mvert,fill=none,line width=1.5pt] (29,29) -- (117.8612,103.5634);
\path[draw=none,fill=mvert] (117.8612,103.5634) circle[radius=3];
\draw[draw=mvert,fill=none,line width=1.5pt] (29,29) -- (104.7018,110.1803);
\draw[draw=mnavy,fill=none,line width=1.5pt] (29,29) -- (91.3052,114.7558);
\path[draw=none,fill=morange] (91.3052,114.7558) circle[radius=3];
\draw[draw=mvert,fill=none,line width=1.5pt] (29,29) -- (77.9658,117.3366);
\draw[draw=mvert,fill=none,line width=1.5pt] (29,29) -- (71.7052,134.699);
\path[draw=none,fill=mvert] (71.7052,134.699) circle[radius=3];
\draw[draw=mnavy,fill=none,line width=1.5pt] (29,29) -- (57.2113,134.2859);
\draw[draw=mvert,fill=none,line width=1.5pt] (29,29) -- (43.474,131.9879);
\path[draw=none,fill=mvert] (43.474,131.9879) circle[radius=3];
\path[draw=none,fill=mnavy] (29,29) circle[radius=10];
\draw[draw=morange,fill=none,line width=2.7pt] (78.8855,39.6035) arc[start angle=12,end angle=87,x radius=51,y radius=51];
\end{scope}
\end{scope}
\end{scope}
}

%% D01 | Contributions coordonnées
\expandafter\def\csname MotifD01\endcsname{%
\begin{scope}
\begin{scope}[shift={(0,0)}]
\begin{scope}[xscale=1,yscale=1]
\foreach \dx in {8,23,...,143} {\foreach \dy in {7,22,...,97} {\fill[mpoint] (\dx,\dy) circle[radius=0.65];}}
\draw[draw=mnavy,fill=none,line width=1.9pt] (124,50) -- (91,50);
\path[draw=morange,fill=mblanc,line width=1.9pt] (124,50) circle[radius=8];
\draw[draw=mnavy,fill=none,line width=1.9pt] (99.5,82.909) -- (83,63.8564);
\path[draw=morange,fill=mblanc,line width=1.9pt] (99.5,82.909) circle[radius=8];
\draw[draw=mnavy,fill=none,line width=1.9pt] (50.5,82.909) -- (67,63.8564);
\path[draw=morange,fill=mblanc,line width=1.9pt] (50.5,82.909) circle[radius=8];
\draw[draw=mnavy,fill=none,line width=1.9pt] (26,50) -- (59,50);
\path[draw=morange,fill=mblanc,line width=1.9pt] (26,50) circle[radius=8];
\draw[draw=mnavy,fill=none,line width=1.9pt] (50.5,17.091) -- (67,36.1436);
\path[draw=morange,fill=mblanc,line width=1.9pt] (50.5,17.091) circle[radius=8];
\draw[draw=mnavy,fill=none,line width=1.9pt] (99.5,17.091) -- (83,36.1436);
\path[draw=morange,fill=mblanc,line width=1.9pt] (99.5,17.091) circle[radius=8];
\path[draw=none,fill=morange] (75,50) circle[radius=19];
\path[draw=none,fill=mblanc] (75,50) circle[radius=9];
\end{scope}
\end{scope}
\end{scope}
}

%% D02 | Arbitrage vers une sortie
\expandafter\def\csname MotifD02\endcsname{%
\begin{scope}
\begin{scope}[shift={(0,0)}]
\begin{scope}[xscale=1,yscale=1]
\foreach \dx in {8,23,...,143} {\foreach \dy in {7,22,...,97} {\fill[mpoint] (\dx,\dy) circle[radius=0.65];}}
\path[draw=mnavy,fill=mblanc,line width=1.7pt] (18,23) circle[radius=7];
\path[draw=morange,fill=none,line width=1.8pt] (26,23) .. controls (55,23) and (56,50) .. (86,50);
\path[draw=mnavy,fill=mnavy,line width=1.7pt] (18,50) circle[radius=7];
\path[draw=mnavy,fill=none,line width=1.8pt] (26,50) .. controls (55,50) and (56,50) .. (86,50);
\path[draw=mnavy,fill=mblanc,line width=1.7pt] (18,78) circle[radius=7];
\path[draw=morange,fill=none,line width=1.8pt] (26,78) .. controls (55,78) and (56,50) .. (86,50);
\path[draw=none,fill=mnavy,rounded corners=2pt] (85,36) rectangle (102,64);
\draw[draw=morange,fill=none,line width=2.8pt] (102,50) -- (119,50);
\draw[draw=morange,fill=none,line width=2.8pt] (119,50) -- (139,50);
\path[draw=none,fill=morange] (139,50) -- (135.2366,51.8646) -- (135.2366,48.1354) -- cycle;
\path[draw=none,fill=mblanc] (92,50) circle[radius=4];
\end{scope}
\end{scope}
\end{scope}
}

%% D03 | Intersection des exigences
\expandafter\def\csname MotifD03\endcsname{%
\begin{scope}
\begin{scope}[shift={(0,0)}]
\begin{scope}[xscale=1,yscale=1]
\foreach \dx in {8,23,...,143} {\foreach \dy in {7,22,...,97} {\fill[mpoint] (\dx,\dy) circle[radius=0.65];}}
\path[draw=mnavy,fill=none,line width=2pt] (59,56) circle[radius=28];
\path[draw=morange,fill=none,line width=2pt] (91,56) circle[radius=28];
\path[draw=mnavy,fill=none,line width=2pt] (75,34) circle[radius=28];
\path[draw=none,fill=morange,rounded corners=1.5pt] (69,48) rectangle (82,61);
\path[draw=none,fill=mblanc] (75,54) circle[radius=4];
\end{scope}
\end{scope}
\end{scope}
}

%% D04 | Alignement progressif
\expandafter\def\csname MotifD04\endcsname{%
\begin{scope}
\begin{scope}[shift={(0,0)}]
\begin{scope}[xscale=1,yscale=1]
\foreach \dx in {8,23,...,143} {\foreach \dy in {7,22,...,97} {\fill[mpoint] (\dx,\dy) circle[radius=0.65];}}
\begin{scope}
\begin{scope}[shift={(22,45)}]
\begin{scope}[rotate=40]
\path[draw=mnavy,fill=mblanc,line width=2pt,rounded corners=1pt] (-5,-17) rectangle (5,17);
\end{scope}
\end{scope}
\end{scope}
\begin{scope}
\begin{scope}[shift={(50,45)}]
\begin{scope}[rotate=29]
\path[draw=mnavy,fill=mblanc,line width=2pt,rounded corners=1pt] (-5,-17) rectangle (5,17);
\end{scope}
\end{scope}
\end{scope}
\begin{scope}
\begin{scope}[shift={(78,45)}]
\begin{scope}[rotate=18]
\path[draw=mnavy,fill=mblanc,line width=2pt,rounded corners=1pt] (-5,-17) rectangle (5,17);
\end{scope}
\end{scope}
\end{scope}
\begin{scope}
\begin{scope}[shift={(106,45)}]
\begin{scope}[rotate=7]
\path[draw=mnavy,fill=mblanc,line width=2pt,rounded corners=1pt] (-5,-17) rectangle (5,17);
\end{scope}
\end{scope}
\end{scope}
\path[draw=none,fill=morange,rounded corners=1pt] (127,30) rectangle (139,70);
\draw[draw=morange,fill=none,line width=1.3pt] (22,16) -- (142,16);
\path[draw=none,fill=morange] (142,16) -- (138.2366,17.8646) -- (138.2366,14.1354) -- cycle;
\end{scope}
\end{scope}
\end{scope}
}

%% D05 | Approche itérative
\expandafter\def\csname MotifD05\endcsname{%
\begin{scope}
\begin{scope}[shift={(0,0)}]
\begin{scope}[xscale=1,yscale=1]
\foreach \dx in {8,23,...,143} {\foreach \dy in {7,22,...,97} {\fill[mpoint] (\dx,\dy) circle[radius=0.65];}}
\path[draw=mnavy,fill=mblanc,line width=1.4pt] (119.1656,66.0749) circle[radius=8];
\path[draw=mnavy,fill=mblanc,line width=1.4pt] (99.1571,85.8144) circle[radius=7.43];
\path[draw=mnavy,fill=mblanc,line width=1.4pt] (73.625,89.376) circle[radius=6.86];
\path[draw=mnavy,fill=mblanc,line width=1.4pt] (53.0825,78.0532) circle[radius=6.29];
\path[draw=mnavy,fill=mblanc,line width=1.4pt] (44.4319,58.7653) circle[radius=5.72];
\path[draw=mnavy,fill=mblanc,line width=1.4pt] (48.6886,40.4234) circle[radius=5.15];
\path[draw=mnavy,fill=mblanc,line width=1.4pt] (61.4675,29.9373) circle[radius=4.58];
\path[draw=mnavy,fill=mblanc,line width=1.4pt] (75.7119,29.6124) circle[radius=4.01];
\path[draw=mnavy,fill=mblanc,line width=1.4pt] (85.22,36.919) circle[radius=3.44];
\path[draw=none,fill=morange] (87.3041,46.4718) circle[radius=2.87];
\path[draw=none,fill=morange] (83.4572,53.0782) circle[radius=2.3];
\path[draw=none,fill=morange] (75,50) circle[radius=7];
\draw[draw=mnavy,fill=none,line width=1.25pt] (119.0166,40.644) arc[start angle=-12,end angle=208,x radius=45,y radius=45];
\end{scope}
\end{scope}
\end{scope}
}

%% D06 | Décision collégiale
\expandafter\def\csname MotifD06\endcsname{%
\begin{scope}
\begin{scope}[shift={(0,0)}]
\begin{scope}[xscale=1,yscale=1]
\foreach \dx in {8,23,...,143} {\foreach \dy in {7,22,...,97} {\fill[mpoint] (\dx,\dy) circle[radius=0.65];}}
\draw[draw=mnavy,fill=none,line width=1.5pt] (35,75) -- (75,50);
\path[draw=mnavy,fill=mblanc,line width=1.6pt,rounded corners=3pt] (25,65) rectangle (45,85);
\path[draw=none,fill=mvert] (35,75) circle[radius=4];
\draw[draw=mnavy,fill=none,line width=1.5pt] (115,75) -- (75,50);
\path[draw=mnavy,fill=mblanc,line width=1.6pt,rounded corners=3pt] (105,65) rectangle (125,85);
\path[draw=none,fill=mvert] (115,75) circle[radius=4];
\draw[draw=mnavy,fill=none,line width=1.5pt] (35,25) -- (75,50);
\path[draw=mnavy,fill=mblanc,line width=1.6pt,rounded corners=3pt] (25,15) rectangle (45,35);
\path[draw=none,fill=morange] (35,25) circle[radius=4];
\draw[draw=mnavy,fill=none,line width=1.5pt] (115,25) -- (75,50);
\path[draw=mnavy,fill=mblanc,line width=1.6pt,rounded corners=3pt] (105,15) rectangle (125,35);
\path[draw=none,fill=morange] (115,25) circle[radius=4];
\path[draw=none,fill=morange] (75,50) circle[radius=16];
\path[draw=none,fill=mblanc,rounded corners=2pt] (68,43) rectangle (82,57);
\draw[draw=mnavy,fill=none,line width=1pt] (105.3767,62.8941) arc[start angle=23,end angle=153,x radius=33,y radius=33];
\draw[draw=mnavy,fill=none,line width=1pt] (44.853,36.5777) arc[start angle=204,end angle=336,x radius=33,y radius=33];
\end{scope}
\end{scope}
\end{scope}
}


% Typographie et cadres de pages
\newcommand{\Texte}[5]{\node[anchor=north west,inner sep=0pt,text=#4,font={\fontsize{#5pt}{#5pt}\selectfont}] at (#1,#2) {#3};}
\newcommand{\MotifCarte}[7]{%
% #1 = Y haut ; #2 = identifiant ; #3 = rang ; #4 titre ; #5 sens ; #6 usage ; #7 macro dessin
  \pgfmathsetmacro{\ybas}{#1-186}
  \fill[mblanc,rounded corners=10pt] (48,\ybas) rectangle (547.28,#1-3);
  \draw[mligne,line width=.6pt,rounded corners=10pt] (48,\ybas) rectangle (547.28,#1-3);
  \fill[mpapier,rounded corners=8pt] (57,\ybas+8) rectangle (309,\ybas+175);
  \begin{scope}[shift={(70,\ybas+16)},scale=1.44]
    #7
  \end{scope}
  \node[anchor=west,inner sep=0pt,text=morange,font={\bfseries\fontsize{9pt}{10pt}\selectfont}] at (327,#1-31) {#2 / #3};
  \node[anchor=west,inner sep=0pt,text=mnavy,font={\displayfont\bfseries\fontsize{16pt}{17pt}\selectfont}] at (327,#1-60) {#4};
  \node[anchor=north west,inner sep=0pt,text width=204pt,align=left,text=mencre,font={\fontsize{9.3pt}{13.7pt}\selectfont}] at (327,#1-74) {#5};
  \draw[mligne,line width=.6pt] (327,\ybas+42)--(527,\ybas+42);
  \node[anchor=west,inner sep=0pt,text=mgris,font={\itshape\fontsize{8.5pt}{10pt}\selectfont}] at (327,\ybas+23) {#6};
}
\newcommand{\DebutPage}{\null\vspace*{-\baselineskip}\begin{tikzpicture}[remember picture,overlay,shift={(current page.south west)},x=1pt,y=1pt]}
\newcommand{\FinPage}{\end{tikzpicture}\clearpage}



%% Carte : motif sur fond papier, rogné au cadre 150 x 100. \IEcard[échelle]{A01}
\newcommand{\IEcard}[2][1.4]{%
  \begin{tikzpicture}[x=1pt,y=1pt,scale=#1]
    \fill[mpapier!70!white] (0,0) rectangle (150,100);
    \begin{scope}\clip (0,0) rectangle (150,100);\Motif{#2}\end{scope}
    \draw[mligne,line width=0.4pt] (0,0) rectangle (150,100);
  \end{tikzpicture}}

%% Noms des motifs
\newcommand{\IEdef}[2]{\expandafter\gdef\csname IEname@#1\endcsname{#2}}
\IEdef{A01}{Noyau et satellites}\IEdef{A02}{Strates emboîtées}\IEdef{A03}{Maillage modulaire}
\IEdef{A04}{Équilibre des contraintes}\IEdef{A05}{Interface commune}\IEdef{A06}{Ossature assemblée}
\IEdef{B01}{Chaîne de transformations}\IEdef{B02}{Relais successifs}\IEdef{B03}{Branches réunifiées}
\IEdef{B04}{Filtre et sélection}\IEdef{B05}{Boucle de rétroaction}\IEdef{B06}{Distribution ramifiée}
\IEdef{C01}{Dispersion autour d'un repère}\IEdef{C02}{Famille de scénarios}\IEdef{C03}{Centre mobile}
\IEdef{C04}{Domaines superposés}\IEdef{C05}{Perturbation locale}\IEdef{C06}{Éventail de possibilités}
\IEdef{D01}{Contributions coordonnées}\IEdef{D02}{Arbitrage vers une sortie}\IEdef{D03}{Intersection des exigences}
\IEdef{D04}{Alignement progressif}\IEdef{D05}{Approche itérative}\IEdef{D06}{Décision collégiale}
\newcommand{\IEname}[1]{\csname IEname@#1\endcsname}

%% ── 3. Ouvertures de partie (ambiance Marque) ───────────────
%% Composition : motif en tiers supérieur, étiquette et titre au centre,
%% filet et provenance du motif en pied. Rien d'autre sur la page.
\newcommand{\IEopener}[3]{% étiquette, titre, motif
  \clearpage\thispagestyle{empty}\papertexture
  \begin{tikzpicture}[remember picture,overlay,shift={(current page.center)}]
    \begin{scope}[shift={(0,6.9cm)},x=1pt,y=1pt,scale=2.55]
      \begin{scope}[shift={(-75,-50)}]
        \begin{scope}\clip (-6,-6) rectangle (156,106);\Motif{#3}\end{scope}
      \end{scope}
    \end{scope}
    \node[align=center,text width=14cm,text=morange] at (0,0.2)
      {\ttfamily\fontsize{10}{12}\selectfont\MakeUppercase{#1}};
    \node[align=center,text width=14.5cm,text=mnavy,execute at begin node={\hyphenpenalty=10000 \exhyphenpenalty=10000 }] at (0,-1.9)
      {\fontsize{30}{35}\selectfont\bfseries #2};
    \draw[mnavy,line width=0.5pt] (-8.2,-10.2) -- (8.2,-10.2);
    \node[anchor=west,text=mgris] at (-8.2,-10.8) {\ttfamily\fontsize{7.5}{9}\selectfont GUIDE DE RÉDACTION \textperiodcentered\ VERSION 7.0};
    \node[anchor=east,text=mgris] at (8.2,-10.8) {\ttfamily\fontsize{7.5}{9}\selectfont MOTIF #3 \textperiodcentered\ \MakeUppercase{\IEname{#3}}};
  \end{tikzpicture}\clearpage}
%% Partie numérotée et partie non numérotée, avec entrée de table des matières
\newcommand{\IEpart}[2]{%
  \refstepcounter{part}\phantomsection
  \addcontentsline{toc}{part}{\protect\numberline{\thepart}#1}%
  \IEopener{Partie \thepart}{#1}{#2}}
\newcommand{\IEpartstar}[2]{%
  \phantomsection\addcontentsline{toc}{part}{#1}%
  \IEopener{Clôture}{#1}{#2}}

%% ── 4. Bibliothèque de l'ambiance Terre ─────────────────────
\tikzset{
  hand/.style={decorate,decoration={random steps,segment length=3.2mm,amplitude=0.55pt},
               line width=1.15pt,draw=noirenc,line join=round,line cap=round},
  handthin/.style={decorate,decoration={random steps,segment length=2.6mm,amplitude=0.4pt},
               line width=0.7pt,draw=noirenc,line cap=round},
}
\newcommand{\chapterpagebegin}[1]{%
  \clearpage\thispagestyle{empty}%
  \begin{tikzpicture}[remember picture,overlay,shift={(current page.center)}]
    \fill[#1] (-10.6,-15) rectangle (10.6,15);
    \pgfmathsetseed{7}
}
\newcommand{\chapterpagetext}[3]{% numéro, titre, couleur du texte
    \node[align=center,text width=13cm,text=#3,execute at begin node={\hyphenpenalty=10000 \exhyphenpenalty=10000 \lefthyphenmin=99 \righthyphenmin=99 }] at (0,0)
      {{\fontsize{14}{17}\selectfont\bfseries\MakeUppercase{Chapitre #1}}\\[10pt]
       {\fontsize{30}{35}\selectfont\bfseries #2}};
  \end{tikzpicture}\newpage}
\newcommand{\handcircle}[4]{% x y rayon couleur
  \fill[#4] (#1+0.09,#2-0.07) circle (#3);
  \pgfmathparse{(#3)<0.95?1:0}\ifdim\pgfmathresult pt>0pt\draw[noirenc,line width=1.0pt] (#1,#2) circle (#3);\else\draw[hand] (#1,#2) circle (#3);\fi}
\newcommand{\handrect}[5]{% x1 y1 x2 y2 couleur
  \fill[#5] (#1+0.10,#2-0.08) rectangle (#3+0.10,#4-0.08);
  \draw[hand] (#1,#2) rectangle (#3,#4);}
\newcommand{\handtri}[7]{% x1 y1 x2 y2 x3 y3 couleur
  \fill[#7] (#1+0.10,#2-0.08) -- (#3+0.10,#4-0.08) -- (#5+0.10,#6-0.08) -- cycle;
  \draw[hand] (#1,#2) -- (#3,#4) -- (#5,#6) -- cycle;}
\newcommand{\handwave}[4]{% x1 x2 y amplitude
  \draw[hand] (#1,#3) .. controls (#1+2,#3+#4) and (#1+4,#3-#4) .. (#1+6,#3)
    .. controls (#1+8,#3+#4) and (#1+10,#3-#4) .. (#1+12,#3)
    .. controls (#1+14,#3+#4) and (#1+16,#3-#4) .. (#1+18,#3);}

%% ── 5. Aides de mise en page ────────────────────────────────
%% Pastille de couleur (fond clair visible grâce au liseré)
\newcommand{\IEsw}[1]{\fcolorbox{grisclair!45}{#1}{\rule{0pt}{6pt}\rule{6pt}{0pt}}}
%% Légende de motif : identifiant, nom, lecture, usage
\newcommand{\IEinfo}[4]{{\footnotesize\textbf{#1}~\textbf{#2.} #3 \textit{Usage~: #4.}}}
%% Miniature d'une page d'un PDF, encadrée d'un filet fin : fichier, page, largeur
\newcommand{\IEthumb}[3]{{\setlength{\fboxsep}{0pt}\setlength{\fboxrule}{0.4pt}\fcolorbox{grisclair!60}{white}{\includegraphics[page=#2,width=#3]{#1}}}}

%% ── 6. Ouvertures de partie, ambiance Terre (guide v7.0) ─────
%% Étiquette « PARTIE n » (ou « CLÔTURE »), titre, couleur du texte
\newcommand{\IEterretext}[3]{%
    \node[align=center,text width=13cm,text=#3,execute at begin node={\hyphenpenalty=10000 \exhyphenpenalty=10000 \lefthyphenmin=99 \righthyphenmin=99 }] at (0,0)
      {{\fontsize{14}{17}\selectfont\bfseries\MakeUppercase{#1}}\\[10pt]
       {\fontsize{30}{35}\selectfont\bfseries #2}};
  \end{tikzpicture}\newpage}
\newcommand{\IEpartT}[2]{%
  \refstepcounter{part}\phantomsection
  \addcontentsline{toc}{part}{\protect\numberline{\thepart}#1}%
  \csname IEterre@#2\endcsname{Partie \thepart}{#1}}
\newcommand{\IEpartTstar}[2]{%
  \phantomsection\addcontentsline{toc}{part}{#1}%
  \csname IEterre@#2\endcsname{Fin du guide}{#1}}

%% 1 · Support physique : feuille, marges, règle, cellules de grille
\expandafter\newcommand\csname IEterre@1\endcsname[2]{\chapterpagebegin{tcotta}
    \handrect{-4.2}{5.2}{3.6}{13.6}{creme}
    \draw[handthin,dashed] (-3.3,6.0) rectangle (2.7,12.8);
    \foreach \y in {11.4,10.4,9.4,8.4}{\draw[handthin] (-2.6,\y) -- (2.0,\y);}
    \draw[hand,line width=1.0pt,creme] (-6.6,5.4) -- (-6.6,13.4);
    \foreach \y in {6.0,7.0,8.0,9.0,10.0,11.0,12.0,13.0}{\draw[handthin,creme] (-6.6,\y) -- (-6.0,\y);}
    \handcircle{7.0}{11.0}{1.3}{beige}
    \handrect{-8.2}{-6.0}{-3.2}{-9.0}{creme}\handrect{-2.6}{-6.0}{2.4}{-9.0}{beige}\handrect{3.0}{-6.0}{8.0}{-9.0}{creme}
    \handrect{-9.8}{-13.4}{9.8}{-12.0}{beige}
  \IEterretext{#1}{#2}{creme}}
%% 2 · Typographie : un « A » géométrique, lignes de texte, caractères
\expandafter\newcommand\csname IEterre@2\endcsname[2]{\chapterpagebegin{beige}
    \handtri{-3.6}{5.4}{0.0}{13.4}{3.6}{5.4}{tcotta}
    \handrect{-1.7}{7.6}{1.7}{8.5}{creme}
    \foreach \y in {12.4,11.2,10.0}{\draw[hand] (5.2,\y) -- (9.4,\y);}
    \handcircle{-7.4}{11.4}{1.2}{sable}
    \handcircle{-6.2}{-8.0}{1.5}{tcotta}\handrect{-3.4}{-9.6}{-0.4}{-6.6}{creme}\handtri{1.2}{-9.6}{2.8}{-6.4}{4.4}{-9.6}{sable}
    \draw[hand] (-8.8,-10.4) -- (8.8,-10.4);
    \handcircle{7.2}{-7.6}{0.6}{tcotta}
  \IEterretext{#1}{#2}{noirenc}}
%% 3 · Structuration : niveaux emboîtés, escalier
\expandafter\newcommand\csname IEterre@3\endcsname[2]{\chapterpagebegin{argile}
    \handrect{-8.0}{12.6}{6.0}{13.6}{tcotta}\handrect{-6.5}{11.0}{4.5}{12.0}{creme}
    \handrect{-5.0}{9.4}{3.0}{10.4}{sable}\handrect{-3.5}{7.8}{1.5}{8.8}{tcotta}
    \handcircle{7.6}{9.6}{1.4}{creme}
    \handrect{-3.2}{-6.0}{3.2}{-12.4}{creme}\handrect{-2.0}{-7.4}{2.0}{-11.0}{tcotta}\handcircle{0}{-9.2}{0.9}{creme}
    \handrect{-9.8}{-13.4}{-5.6}{-12.4}{sable}
  \IEterretext{#1}{#2}{noirenc}}
%% 4 · Pagination : pile de pages, suite de pages numérotées
\expandafter\newcommand\csname IEterre@4\endcsname[2]{\chapterpagebegin{sable}
    \handrect{-4.8}{5.6}{1.2}{13.2}{beige}\handrect{-4.0}{5.1}{2.0}{12.7}{creme}\handrect{-3.2}{4.6}{2.8}{12.2}{beige}
    \foreach \y in {11.0,10.0,9.0}{\draw[handthin] (-2.4,\y) -- (1.8,\y);}
    \handcircle{7.2}{10.0}{1.5}{tcotta}
    \foreach \x/\c in {-7/creme,-3.5/tcotta,0/creme,3.5/tcotta,7/creme}{\handrect{\x-1.0}{-7.0}{\x+1.0}{-9.0}{\c}}
    \foreach \x in {-5.2,-1.7,1.8,5.3}{\draw[hand,-{Stealth[length=2.4mm]}] (\x,-8.0) -- (\x+0.9,-8.0);}
  \IEterretext{#1}{#2}{noirenc}}
%% 5 · En-tête et pied de page : feuille à deux bandeaux
\expandafter\newcommand\csname IEterre@5\endcsname[2]{\chapterpagebegin{creme}
    \handrect{-4.6}{5.0}{4.6}{13.6}{beige}
    \handrect{-4.6}{12.0}{4.6}{13.6}{tcotta}\handrect{-4.6}{5.0}{4.6}{6.4}{tcotta}
    \foreach \y in {10.6,9.6,8.6,7.6}{\draw[handthin] (-3.4,\y) -- (3.4,\y);}
    \handcircle{8.0}{11.0}{1.0}{tcottaL}
    \handrect{-9.0}{-6.8}{9.0}{-8.0}{tcotta}\handrect{-9.0}{-10.0}{9.0}{-11.2}{beige}
    \handcircle{-7.2}{-9.0}{0.5}{tcotta}\handcircle{7.2}{-9.0}{0.5}{tcotta}
  \IEterretext{#1}{#2}{noirenc}}
%% 6 · Éléments visuels : histogramme, tableau, courbe
\expandafter\newcommand\csname IEterre@6\endcsname[2]{\chapterpagebegin{tcottaL}
    \handrect{-8.0}{5.6}{-6.6}{8.6}{creme}\handrect{-6.0}{5.6}{-4.6}{11.0}{tcotta}\handrect{-4.0}{5.6}{-2.6}{9.6}{creme}\handrect{-2.0}{5.6}{-0.6}{12.4}{beige}
    \draw[hand] (-8.6,5.6) -- (0.2,5.6);
    \handrect{2.2}{6.0}{8.6}{12.4}{creme}
    \foreach \y in {8.1,10.2}{\draw[handthin] (2.2,\y) -- (8.6,\y);}
    \foreach \x in {4.3,6.4}{\draw[handthin] (\x,6.0) -- (\x,12.4);}
    \handwave{-9.0}{0}{-7.6}{0.7}\handwave{-9.0}{0}{-9.2}{0.7}
    \handcircle{6.4}{-8.0}{1.6}{creme}\handtri{1.0}{-10.6}{2.8}{-6.6}{4.6}{-10.6}{beige}
  \IEterretext{#1}{#2}{noirenc}}
%% 7 · Style éditorial : les cinq formes élémentaires et l'ombre décalée
\expandafter\newcommand\csname IEterre@7\endcsname[2]{\chapterpagebegin{tcotta}
    \handcircle{-7.0}{9.4}{1.3}{creme}\handrect{-4.6}{8.2}{-2.2}{10.6}{beige}\handtri{-1.4}{8.2}{0.2}{10.9}{1.8}{8.2}{sable}
    \draw[hand,-{Stealth[length=3mm]}] (3.4,9.4) -- (6.4,9.4);\draw[hand] (7.2,9.4) -- (9.0,9.4);
    \handwave{-9.0}{0}{12.4}{0.5}
    \handrect{-3.4}{-6.0}{3.4}{-12.4}{creme}
    \draw[handthin,dashed] (-3.4,-6.0) rectangle (3.4,-12.4);
    \handcircle{7.0}{-9.0}{1.2}{beige}\handcircle{-7.0}{-8.0}{0.6}{sable}
  \IEterretext{#1}{#2}{creme}}
%% 8 · Normes de rédaction et de citation : livre ouvert, guillemets, crochets
\expandafter\newcommand\csname IEterre@8\endcsname[2]{\chapterpagebegin{argile}
    \handrect{-6.2}{6.0}{-0.2}{12.6}{creme}\handrect{0.2}{6.0}{6.2}{12.6}{beige}
    \foreach \y in {11.2,10.2,9.2,8.2}{\draw[handthin] (-5.2,\y) -- (-1.2,\y);\draw[handthin] (1.2,\y) -- (5.2,\y);}
    \handcircle{-2.4}{-8.0}{1.5}{tcotta}\handtri{-3.6}{-9.2}{-1.6}{-9.2}{-2.6}{-11.2}{tcotta}
    \handcircle{2.4}{-8.0}{1.5}{tcotta}\handtri{1.2}{-9.2}{3.2}{-9.2}{2.2}{-11.2}{tcotta}
    \draw[hand] (-8.6,-6.0) -- (-9.2,-6.0) -- (-9.2,-11.0) -- (-8.6,-11.0);\draw[hand] (8.6,-6.0) -- (9.2,-6.0) -- (9.2,-11.0) -- (8.6,-11.0);
  \IEterretext{#1}{#2}{noirenc}}
%% 9 · Métadonnées, gouvernance, preuve : sceau et chaîne de preuve
\expandafter\newcommand\csname IEterre@9\endcsname[2]{\chapterpagebegin{beige}
    \handcircle{0}{9.4}{3.4}{tcotta}\handcircle{0}{9.4}{2.2}{creme}\handcircle{0}{9.4}{0.9}{tcotta}
    \handcircle{-7.6}{12.0}{0.8}{sable}\handcircle{7.6}{7.0}{1.1}{argile}
    \handrect{-8.8}{-6.6}{-6.0}{-9.4}{creme}\handrect{-4.4}{-6.6}{-1.6}{-9.4}{tcotta}\handrect{1.6}{-6.6}{4.4}{-9.4}{creme}\handrect{6.0}{-6.6}{8.8}{-9.4}{tcotta}
    \foreach \x in {-6.0,-1.6,4.4}{\draw[hand,line width=1.4pt] (\x,-8.0) -- (\x+1.6,-8.0);}
    \draw[hand] (-9.4,-12.6) -- (9.4,-12.6);
  \IEterretext{#1}{#2}{noirenc}}
%% 10 · Gabarits : une feuille remplie, deux gabarits vides en pointillé
\expandafter\newcommand\csname IEterre@10\endcsname[2]{\chapterpagebegin{sable}
    \handrect{-8.2}{5.4}{-3.0}{13.4}{creme}
    \foreach \y in {12.0,11.0,10.0}{\draw[handthin] (-7.4,\y) -- (-3.8,\y);}
    \draw[handthin,dashed] (-2.2,5.4) rectangle (3.0,13.4);\draw[handthin,dashed] (3.8,5.4) rectangle (9.0,13.4);
    \handcircle{0.4}{9.4}{0.8}{tcotta}
    \handrect{-3.0}{-6.4}{3.0}{-12.0}{tcotta}\handcircle{0}{-9.2}{1.7}{creme}
    \handrect{-9.8}{-13.4}{-6.0}{-12.4}{beige}
  \IEterretext{#1}{#2}{noirenc}}
%% 11 · Production et publication : chaîne de fabrication
\expandafter\newcommand\csname IEterre@11\endcsname[2]{\chapterpagebegin{creme}
    \draw[hand,line width=1.4pt] (-9.0,9.0) -- (9.0,9.0);
    \handrect{-8.0}{9.0}{-5.2}{11.4}{tcotta}\handrect{-3.4}{9.0}{-0.6}{11.4}{beige}\handrect{1.2}{9.0}{4.0}{11.4}{tcottaL}
    \handcircle{7.0}{10.4}{1.4}{tcotta}
    \draw[hand,-{Stealth[length=3mm]}] (-8.0,6.6) -- (6.0,6.6);
    \handcircle{-4.0}{-8.0}{2.2}{beige}\handcircle{-4.0}{-8.0}{0.8}{tcotta}
    \handtri{2.0}{-10.4}{2.0}{-5.6}{6.4}{-8.0}{tcotta}
    \handrect{-9.8}{-13.4}{9.8}{-12.4}{beige}
  \IEterretext{#1}{#2}{noirenc}}
%% 12 · Références du dépôt : arborescence de dossiers
\expandafter\newcommand\csname IEterre@12\endcsname[2]{\chapterpagebegin{tcottaL}
    \handrect{-1.6}{12.0}{1.6}{13.4}{creme}
    \draw[hand] (0,12.0) -- (0,11.0);\draw[hand] (-5.5,11.0) -- (5.5,11.0);
    \foreach \x in {-5.5,0,5.5}{\draw[hand] (\x,11.0) -- (\x,10.0);\handrect{\x-1.5}{8.6}{\x+1.5}{10.0}{beige}}
    \foreach \x in {-5.5,5.5}{\draw[hand] (\x,8.6) -- (\x,7.6);\handrect{\x-1.0}{6.2}{\x+1.0}{7.6}{creme}}
    \handcircle{-6.0}{-8.0}{1.2}{creme}\handcircle{0}{-9.4}{1.6}{tcotta}\handcircle{6.0}{-8.0}{1.2}{beige}
    \draw[hand] (-6.0,-8.0) -- (0,-9.4) -- (6.0,-8.0);
  \IEterretext{#1}{#2}{noirenc}}
%% 13 · Cas particuliers : une forme qui sort de la rangée
\expandafter\newcommand\csname IEterre@13\endcsname[2]{\chapterpagebegin{argile}
    \foreach \x in {-8.0,-4.2,-0.4,7.2}{\handrect{\x-1.2}{9.0}{\x+1.2}{11.4}{creme}}
    \handtri{2.2}{8.8}{3.4}{12.2}{4.6}{8.8}{tcotta}
    \draw[hand,line width=1.4pt] (-9.0,6.8) -- (9.0,6.8);
    \draw[hand,line width=1.2pt] (-9.0,-8.0) -- (-3.0,-8.0) .. controls (-1.0,-8.0) and (0.0,-5.4) .. (2.0,-5.4) .. controls (4.0,-5.4) and (4.6,-8.0) .. (6.0,-8.0) -- (9.0,-8.0);
    \handcircle{2.0}{-5.4}{0.7}{tcotta}\handcircle{-5.4}{-11.0}{1.3}{sable}
  \IEterretext{#1}{#2}{noirenc}}
%% 14 · Annexes : feuilles superposées à onglets
\expandafter\newcommand\csname IEterre@14\endcsname[2]{\chapterpagebegin{beige}
    \handrect{-5.0}{5.6}{3.4}{13.2}{sable}\handrect{-4.2}{5.1}{4.2}{12.7}{creme}\handrect{-3.4}{4.6}{5.0}{12.2}{tcottaL}
    \handrect{5.0}{11.0}{6.4}{12.2}{tcotta}\handrect{5.0}{9.0}{6.4}{10.2}{creme}\handrect{5.0}{7.0}{6.4}{8.2}{tcotta}
    \foreach \y in {10.8,9.8,8.8}{\draw[handthin] (-2.4,\y) -- (3.0,\y);}
    \handcircle{-8.2}{11.0}{1.0}{tcotta}
    \foreach \x/\c in {-6.0/tcotta,-1.0/creme,4.0/sable}{\handrect{\x-1.4}{-6.4}{\x+1.4}{-9.4}{\c}}
    \handrect{-9.8}{-13.4}{9.8}{-12.4}{sable}
  \IEterretext{#1}{#2}{noirenc}}
%% 15 · Glossaire et index : colonnes d'entrées, onglets alphabétiques
\expandafter\newcommand\csname IEterre@15\endcsname[2]{\chapterpagebegin{sable}
    \foreach \y/\l in {12.6/3.6,11.4/5.0,10.2/2.8,9.0/4.4,7.8/3.2,6.6/4.0}{\handcircle{-8.0}{\y}{0.32}{tcotta}\handrect{-7.2}{\y-0.22}{-7.2+\l}{\y+0.22}{creme}}
    \foreach \y/\l in {12.6/4.2,11.4/2.6,10.2/4.8,9.0/3.4,7.8/5.0,6.6/2.4}{\handcircle{1.2}{\y}{0.32}{tcotta}\handrect{2.0}{\y-0.22}{2.0+\l}{\y+0.22}{beige}}
    \foreach \x/\h in {-8.0/2.6,-6.2/3.4,-4.4/2.2,-2.6/3.8,-0.8/2.8,1.0/3.2,2.8/2.0,4.6/3.6,6.4/2.4,8.2/3.0}{\handrect{\x-0.7}{-5.6}{\x+0.7}{-5.6-\h}{tcotta}}
    \handrect{-9.8}{-13.4}{9.8}{-12.4}{beige}
  \IEterretext{#1}{#2}{noirenc}}
%% 16 · Clôture : soleil bas sur l'horizon
\expandafter\newcommand\csname IEterre@16\endcsname[2]{\chapterpagebegin{tcotta}
    \handcircle{0}{9.4}{3.4}{creme}
    \foreach \a in {20,50,80,110,140,170}{\draw[handthin,creme] ({4.2*cos(\a)},{9.4+4.2*sin(\a)}) -- ({5.2*cos(\a)},{9.4+5.2*sin(\a)});}
    \fill[tcotta] (-10.6,5.2) rectangle (10.6,7.0);
    \handwave{-9.0}{0}{6.2}{0.6}\handwave{-9.0}{0}{4.8}{0.6}
    \handrect{-9.8}{-13.4}{9.8}{-12.0}{beige}
    \handrect{-4.0}{-9.0}{4.0}{-6.0}{beige}\handcircle{7.4}{-8.4}{1.1}{creme}\handcircle{-7.4}{-8.4}{0.7}{creme}
  \IEterretext{#1}{#2}{creme}}
